%%%%%%%%%%%%%%%%%%%%%%%%%%%%%%%%%%%%%%%%%%%%%%%%%%%%%%%%%%%%%%%%%%%%%%
%% Shared pieces for the l05_sc sampling spectra.
%%
%% l5_sh, l5_shaaf and l5_subsample are the same drawing three times: no
%% filter, an anti-alias low pass, and a bandpass around f_s. The slides make
%% their point by looking alike, so the axis, the wanted signal and the tones
%% all come from here and every figure inherits the same geometry.
%%
%% Every name is prefixed. \th is a LaTeX built-in and redefining it is a
%% fatal error, so short generic names are not safe.
%%%%%%%%%%%%%%%%%%%%%%%%%%%%%%%%%%%%%%%%%%%%%%%%%%%%%%%%%%%%%%%%%%%%%%

\definecolor{specgreen}{rgb}{0.00,0.50,0.00}  % = armygreen
\definecolor{tonered}{rgb}{1.00,0.00,0.00}   % house red
\definecolor{tonemag}{rgb}{0.00,0.00,1.00}   % house blue; the second tone pair
\definecolor{bandorange}{rgb}{0.00,0.00,0.00}  % filters are annotation: black
\newcommand{\specHalf}{1.2}      % f_s/2
\newcommand{\specXmax}{3.5}      % axis half length
\newcommand{\specVh}{1.9}        % vertical axis height
\newcommand{\specToneH}{1.05}    % tone height
\newcommand{\specSkirt}{0.20}   % filter roll-off width

% Panel baselines, so the two rows line up across all three figures.
\newcommand{\specYtop}{3.9}
\newcommand{\specYbot}{0}

% One frequency axis with its Nyquist ticks. #1 is the baseline.
\newcommand{\specAxis}[1]{%
  \draw[thick, <->] (-\specXmax,#1) -- (\specXmax,#1);
  \draw[thick, ->] (0,#1) -- (0,#1+\specVh);
  \foreach \p/\lbl in {-2/{-f_{s}}, -1/{-f_{s}/2}, 1/{f_{s}/2}, 2/{f_{s}}} {
    \draw[thick] (\p*\specHalf,#1+0.13) -- (\p*\specHalf,#1-0.13);
    \node[anchor=north] at (\p*\specHalf,#1-0.22) {$\lbl$};
  }
}

% The wanted signal, centred on DC. #1 is the baseline.
\newcommand{\specBump}[1]{%
  \draw[specgreen, very thick] plot[smooth, tension=0.7] coordinates
    {(-1.05,#1) (-0.75,#1+0.80) (-0.35,#1+1.30) (0,#1+1.45)
     (0.35,#1+1.30) (0.75,#1+0.80) (1.05,#1)};
}

% A tone: colour, position, baseline.
\newcommand{\specTone}[3]{\draw[#1, very thick, ->] (#2,#3) -- (#2,#3+\specToneH);}

% A filter response, drawn as a trapezoid: centre, half width, height,
% baseline. Used for both the anti-alias low pass and the subsampling
% bandpass, so the two read as the same kind of object.
\newcommand{\specBand}[4]{%
  \draw[bandorange, very thick, rounded corners=3pt]
    (#1-#2-\specSkirt,#4) -- (#1-#2,#4+#3)
    -- (#1+#2,#4+#3) -- (#1+#2+\specSkirt,#4);
}

% A block in the signal chain: left edge x, centre y, label.
\newcommand{\specBlock}[3]{%
  \draw[thick] (#1,#2-0.5) rectangle ++(1.5,1.0);
  \node at (#1+0.75,#2) {#3};
}

% An arrow between chain elements: from x, to x, at y.
\newcommand{\specWire}[3]{\draw[thick, ->] (#1,#3) -- (#2,#3);}
